El control del acceso es una de las tareas diarias de cualquier institución
educativa: decide quién puede entrar al plantel y, cuando se hace bien, deja
constancia de quién entró y a qué hora. En el Tecnológico de Estudios
Superiores de Chalco (TESCHA) esta tarea recae en el personal de vigilancia,
que revisa a simple vista la credencial física de cada estudiante. El método
es sencillo, pero se vuelve lento cuando casi todos llegan al mismo tiempo, no
comprueba que quien muestra la credencial sea su titular y no deja ningún
registro que la institución pueda consultar después.

El proyecto \textit{Escanea y Entra: Prototipo de Acceso QR para el TESCHA}
propone probar otra forma de hacerlo. Cada estudiante muestra en su celular un
código QR que cambia cada pocos segundos; un lector colocado en la entrada lo
valida, abre un torniquete a escala y guarda el ingreso. La prueba se hace en
un solo acceso peatonal, durante el turno vespertino, y compara el acceso
actual con el acceso durante el piloto para saber si la propuesta realmente
ayuda.

Este documento presenta el proyecto desde la gestión de proyectos. Primero
plantea el problema que le da origen. Después explica cómo se inició el
proyecto, qué objetivos persigue, qué se va a construir y hasta dónde llega,
cómo se planificaron el tiempo, el costo y los riesgos, quién hace cada parte,
cómo se cuidará la calidad, cómo se dará seguimiento al avance y cómo se
medirán los resultados. Cierra con el estado actual y los próximos pasos. Para
ello reúne lo que el equipo ha elaborado hasta ahora: la ficha inicial, el
plan de calidad, el tablero de metas, la presentación del avance y el
documento de investigación. Las partes técnicas se explican con palabras
sencillas.

Algunos datos todavía no existen porque el proyecto apenas arranca. Cuando
una cifra depende de las mediciones del diagnóstico o del piloto, el documento
la deja como pendiente y no la estima.
